% project-id: q-a-rag
% project-name: RAG Document Question Answering System
% repository-url: https://github.com/kartikaysaxena/q-a-rag
% source-revision: e94ddfa97d5a0d6841e4a02b84faf10c5347d1fd
% role-families: machine learning, ai engineering, rag, python
% skills: Python, RAG, LangChain, FAISS, Streamlit, Google Generative AI Embeddings, Groq, Llama 3, PDF Processing
\ifdefined\ResumeProjectMode
  \def\ProjectDocumentStart{}
  \def\ProjectDocumentEnd{}
\else
  \documentclass[11pt,border=12pt,varwidth=200mm]{standalone}
  \NeedsTeXFormat{LaTeX2e}
\ProvidesPackage{project-blob}[2026/09/20 Reusable resume project block]

\RequirePackage[T1]{fontenc}
\RequirePackage{lmodern}
\RequirePackage{xcolor}
\RequirePackage{enumitem}
\RequirePackage{hyperref}

\definecolor{ProjectAccent}{HTML}{1F4E79}
\definecolor{ProjectMuted}{HTML}{4B5563}
\hypersetup{colorlinks=true,urlcolor=ProjectAccent}

% #1 is a stable bullet ID, #2 is a comma-separated JD matching tag list,
% and #3 is the rendered resume bullet. The first two arguments are kept in
% source so a resume builder can select individual bullets deterministically.
\newcommand{\ProjectBullet}[3]{\item #3}

% title, repository URL, role families, skills, summary, bullet commands
\newcommand{\ProjectBlob}[6]{%
  \noindent\begin{minipage}{\linewidth}
    \raggedright
    {\Large\bfseries\color{ProjectAccent} #1}\hfill
    {\small\href{#2}{Repository}}\par
    \vspace{2pt}
    {\small\color{ProjectMuted}\textbf{Role fit:} #3}\par
    {\small\color{ProjectMuted}\textbf{Skills:} #4}\par
    \vspace{5pt}
    #5\par
    \vspace{3pt}
    \begin{itemize}[leftmargin=1.35em,itemsep=2pt,topsep=2pt,parsep=0pt]
      #6
    \end{itemize}
  \end{minipage}%
}

  \def\ProjectDocumentStart{\begin{document}\resumeSubHeadingListStart}
  \def\ProjectDocumentEnd{\resumeSubHeadingListEnd\end{document}}
\fi
\ProjectDocumentStart

\resumeSubheading
{ RAG Document Question Answering System} %Project Name
{\href{https://github.com/kartikaysaxena/q-a-rag}{Github}}
{A Python RAG system for answering questions from PDF document collections}
{}

\resumeItemListStart
% bullet-id: document-qa-app
% bullet-skills: Python, Streamlit, RAG, PDF Processing
\item {Built a Python/Streamlit RAG application for grounded question answering across local PDF collections through LangChain.}
% bullet-id: retrieval-pipeline
% bullet-skills: LangChain, FAISS, Google Generative AI Embeddings, Text Chunking
\item {Chunked PDFs into 1,000-character windows with 200-character overlap and indexed Google embeddings in FAISS locally.}
% bullet-id: grounded-generation
% bullet-skills: Retrieval, LangChain, Groq, Llama 3, Prompt Engineering
\item {Grounded Groq-hosted Llama 3 responses in FAISS-retrieved document context through a retrieval chain for factual answers.}
% bullet-id: interactive-workflow
% bullet-skills: Streamlit, State Management, Performance Measurement, User Experience
\item {Cached document embeddings in Streamlit session state to accelerate repeated queries and inspect in-app response latency.}
\resumeItemListEnd
\vspace{-2mm}

\ProjectDocumentEnd
